\section{Introduction}

\subsection{The Computational Demands of Cislunar Mission Planning}

With the Artemis program, the Gateway station, the Chang'e lunar exploration
program, and commercial lunar missions advancing, cislunar space has become a
focal region for mission design. Unlike near-Earth orbits, cislunar dynamics
are dominated by the combined gravity of the Earth and the Moon: periodic
orbits near the libration points (Halo, NRHO, DRO, and others) have no
analytic solution, and their design relies on numerical methods in the
three-body framework---differential correction, multiple shooting, parameter
continuation, and invariant manifold theory---while transfer design between
orbits must trade off impulsive, low-energy, and low-thrust propulsion modes.
Such computations share three traits: \textbf{many models} (the circular
restricted three-body problem CR3BP, ephemeris $N$-body, and the bicircular
restricted four-body problem BCR4BP coexist at different fidelities),
\textbf{heavy iteration} (a single ephemeris correction often involves tens
of thousands of trajectory segments), and \textbf{long chains} (initial
guess, correction, high-fidelity propagation, and control simulation are
tightly coupled).

\subsection{Positioning: A Numerical Foundation for the LLM+Agent Era}

e2m2e (Earth to Moon, Moon to Earth) addresses these needs as an
\textbf{algorithm toolset infrastructure for cislunar mission planning}. In
an LLM+Agent-style autonomous mission planning system, the large language
model understands mission intent and decomposes/orchestrates subtasks; e2m2e
handles the numerical half---building dynamical models of cislunar space,
generating periodic orbit families, designing transfer paths between orbits,
and visualizing results for inspection. Beyond the conventional Python API
and CLI, e2m2e natively provides an MCP (Model Context Protocol) server: 13
task-level tools are exposed to LLM agents over stdio, the tool list is
derived from Facade method metadata, artifacts are automatically archived,
and \texttt{record\_id}s chain across tools---so that ``design an NRHO and
run 100 Monte Carlo station-keeping simulations on it'' becomes a single
executable instruction in natural language.

\subsection{Relationship to Existing Tools}

\begin{itemize}
  \item \textbf{GMAT / STK}: general-purpose mission analysis tools, driven
        mainly by GUIs and scripts, with thorough validation data but hardly
        callable programmatically by agents;
  \item \textbf{Orekit / poliastro}: general-purpose astrodynamics libraries
        covering near-Earth to deep space, but not centered on the cislunar
        three-body problem and offering no task-level encapsulation;
  \item \textbf{e2m2e}: focused on cislunar space, encapsulating the entire
        design chain---family selection, initial guess, correction,
        high-fidelity propagation, transfer, station keeping---as task-level
        interfaces, with a Rust numerical kernel and an agent-native MCP
        interface.
\end{itemize}

e2m2e does not aim to replace general-purpose tools but complements them:
propagation accuracy is aligned with GMAT and DFH to the sub-100\,m level
(Section~\ref{sec:verification}), the coordinate and time chains are
cross-validated against GMAT reference data, and it offers what general tools
lack---periodic orbit family generation and an agent-native interface.
Section~\ref{sec:comparison} develops a systematic comparison against ten
peer and ecosystem tools based on a code-level survey.

\subsection{Organization of This Report}

Section~\ref{sec:architecture} presents the system architecture: five
architectural modules and a layered ``Python orchestration + Rust numerical
kernel'' runtime. Section~\ref{sec:dynamics} covers the spacetime reference,
dynamical models, and high-fidelity force models.
Sections~\ref{sec:orbit} and~\ref{sec:transfer} present the core algorithms
for periodic orbit design and transfer design. Section~\ref{sec:control}
covers orbit control and Monte Carlo simulation. Section~\ref{sec:interfaces}
describes the Facade/CLI/MCP interfaces and the orbit catalog.
Section~\ref{sec:comparison} positions e2m2e against ten tools including
GMAT, STK, ATK, Orekit, and pykep.
Section~\ref{sec:verification} presents the verification framework and
concludes with an outlook.
